\chapter{Von Neumann--Bernays--Gödel set theory}
Throughout the exposition, we have been working in the von Neumann--Bernays--Gödel ($\mathsf{NBG}$) set theory, as Part I of \cite{LambekScott1994} does.

Let us recall the main characteristics of the $\mathsf{NBG}$ set theory and its conservativity result over $\mathsf{ZFC}$.

\section{Sets and classes}
The language of $\mathsf{NBG}$ has two sorts of variables: lower-case letters $x,y,z,\ldots$ symbolize sets, while upper-case letters $X,Y,Z,\ldots$ denote classes, and there's a binary relation symbol $\in$ for membership.
Only sets can be members of classes. Every set $a$ determines the class $\{x : x\in a\}$, and we identify the set with this class. A class which cannot be identified with a set in this way is called a \emph{proper class}.
We denote by
\[ V := \{x : x = x\} \]
the class of all sets.

A \emph{class function} is a class $F$ of ordered pairs such that every set in its domain has a unique image; if $(x,y)\in F$, we write $F(x) = y$.

We use $\mathsf{NBG}$ with global choice. In particular, ordinary set-theoretic choice holds. Nevertheless, we list it separately below in order to keep the familiar $\mathsf{ZFC}$ axioms visible.
The need for global choice will be discussed below in the section concerning the sizes of the categories we use.

\section{Axioms}
We use the following convenient schematic presentation of $\mathsf{NBG}$. It keeps the $\mathsf{ZFC}$ axioms explicit and adds the axioms specific to $\mathsf{NBG}$.
The list is not minimal, in the sense that some axioms can be derived from others.

\begin{itemize}
    \item[(\textnormal{NBG1})] \textbf{Axiom of the empty set.}
    There exists an empty set:
    \[\exists x\;\forall y\;y\notin x.\]
    \item[(\textnormal{NBG2})] \textbf{Axiom of extensionality for sets.}
    A set is determined by the collection of its elements, i.e., sets coincide if and only if they have the same elements:
    \[\forall a\;\forall b\;\left(a = b \leftrightarrow \forall x\,(x\in a \leftrightarrow x\in b)\right).\]
    \item[(\textnormal{NBG3})] \textbf{Axiom schema of separation.}
    If $a$ is a set and $\varphi(x)$ is a formula in the language of set theory, then $\{x\in a : \varphi(x)\}$ is a set:
    \[\forall a\;\exists b\;\forall x\;\left(x\in b \leftrightarrow (x\in a\land\varphi(x))\right).\]
    Here $b$ does not occur freely in $\varphi$.
    \item[(\textnormal{NBG4})] \textbf{Axiom of pairing.}
    Given sets $a$ and $b$, there exists the set $\{a,b\}$:
    \[\forall a\;\forall b\;\exists p\;\forall x\;\left(x\in p \leftrightarrow (x = a\lor x = b)\right).\]
    \item[(\textnormal{NBG5})] \textbf{Axiom of union.}
    Given a set $a$, there exists a set whose elements are precisely the elements of the elements of $a$:
    \[\forall a\;\exists b\;\forall x\;\left(x\in b \leftrightarrow \exists y\in a\;x\in y\right).\]
    \item[(\textnormal{NBG6})] \textbf{Power set axiom.}
    Given a set $a$, there exists the set $\ps(a)$ whose elements are the subsets of $a$:
    \[\forall a\;\exists b\;\forall x\;\left(x\in b \leftrightarrow x\subseteq a\right).\]
    \item[(\textnormal{NBG7})] \textbf{Axiom of infinity.}
    There exists an inductive set:
    \[\exists a\;\left(\emptyset\in a\land\forall x\in a\;(x\cup\{x\}\in a)\right).\]
    \item[(\textnormal{NBG8})] \textbf{Axiom of replacement.}
    If $a$ is a set and $F : V \to V$ is a class function, then
    \[F[a] := \{F(x) : x\in a\}\]
    is a set. In symbols:
    \[\forall F\;\left((F : V \to V) \rightarrow \forall a\;\exists b\;\forall y\;(y\in b \leftrightarrow \exists x\in a\;y = F(x))\right).\]
    \item[(\textnormal{NBG9})] \textbf{Axiom of choice.}
    Given a set $a$ of nonempty sets, there exists a choice function $f : a \to \bigcup a$ such that $f(x)\in x$ for every $x\in a$:
    \[\forall a\;\left(\emptyset\notin a \rightarrow \exists f\;(f : a \to \bigcup a\land\forall x\in a\;f(x)\in x)\right).\]
    \item[(\textnormal{NBG10})] \textbf{Axiom of foundation.}
    Every nonempty set has an $\in$-minimal element:
    \[\forall a\;\left(a\neq\emptyset \rightarrow \exists x\in a\;(x\cap a = \emptyset)\right).\]
    \item[(\textnormal{NBG11})] \textbf{Axiom of extensionality for classes.}
    A class is determined by the collection of its members, i.e., two classes coincide if and only if they have the same members:
    \[\forall X\;\forall Y\;\left(X = Y \leftrightarrow \forall x\;(x\in X \leftrightarrow x\in Y)\right).\]
    \item[(\textnormal{NBG12})] \textbf{Axiom schema of elementary class comprehension.}
    If $\varphi(x)$ is a formula in the language of $\mathsf{NBG}$, with free variable $x$, and which does not contain quantifiers over classes, then there exists the class of all sets which satisfy $\varphi$:
    \[\exists X\;\forall x\;\left(x\in X \leftrightarrow \varphi(x)\right).\]
    Here $X$ does not occur freely in $\varphi$.
    \item[(\textnormal{NBG13})] \textbf{Axiom of global choice.}
    There exists a class function which chooses an element from every nonempty set:
    \[\exists F\;\left(F : V \to V \land\forall x\;(x\neq\emptyset \rightarrow F(x)\in x)\right).\]
\end{itemize}

The list of axioms roughly follows the exposition given in \cite[\S\S~2.1--2.2, pp.~9--19]{Schindler2014}.

\section{Conservativity over \texorpdfstring{$\mathsf{ZFC}$}{ZFC}}

\begin{theorem}[Conservativity of $\mathsf{NBG}$ over $\mathsf{ZFC}$]\label{thm:nbg-conservativity}
    For every sentence $\varphi$ in the language of set theory, and hence containing no class variables,
    \[\mathsf{NBG}\vdash\varphi \Longleftrightarrow \mathsf{ZFC}\vdash\varphi.\]
    In other words, every theorem about sets which can be proved using classes can already be proved in $\mathsf{ZFC}$.
\end{theorem}

See \cite[Thm.~1, pp.~44--59]{Felgner1971} for a proof of this theorem.

\section{Sizes of the categories used}

\begin{definition}[Small category]
    A category $\mathcal C$ is \emph{small} if both $\ob(\mathcal C)$ and $\mor(\mathcal C)$ are sets in the sense of $\mathsf{NBG}$.
\end{definition}

\begin{definition}[Locally small category]
    We say that a category $\mathcal C$ is \emph{locally small} if $\Hom_{\mathcal C}(A,B)$ is a set, in the sense of $\mathsf{NBG}$, for every pair of objects $A,B$.
\end{definition}

\begin{remark}[Smallness implies local smallness]
    Every small category is locally small, but the converse is not necessarily true.
\end{remark}

\begin{example}
    The category $\mathsf{Set}$ is locally small and not small.
\end{example}

\begin{definition}[Small calculus]
    A typed $\lambda$-calculus $\mathcal L$ is \emph{small} if both $\mathsf{Ty}(\mathcal L)$ and $\mathsf{Tm}(\mathcal L)$ are sets in the sense of $\mathsf{NBG}$.
\end{definition}

For all three versions of the Curry--Howard--Lambek correspondence, we take the objects on the categorical side to be small categories and those on the syntactic side to be small calculi.
The resulting categories of structures are large but locally small. Under this convention, the constructions of polynomial categories, $\mathbf L(\mathcal C)$ and $\mathbf C(\mathcal L)$ preserve smallness, so the three correspondences are ordinary equivalences in $\mathsf{NBG}$.
Local smallness alone would not suffice: a category with a proper class of objects cannot itself be an object of another category in $\mathsf{NBG}$.

The ambient categories in the fixed-point applications need not be small. If such a category is large but locally small, the Curry--Howard--Lambek argument can be carried out in the small full subcategory generated by the objects involved and closed under the terminal object,
products and, when required, exponentials and the chosen WNNO. Local smallness makes its collection of arrows a set. Lawvere's theorem itself requires no size assumption.

Finally, the ordinary axiom of choice suffices for non-unique iterator choices in a small category; by contrast, global choice suffices to make such choices uniformly over a proper class of objects,
when all the iterators have to be fixed simultaneously. In the NNO case, the uniqueness removes the need for choice.
